\section{Introduction}
\label{sec:intro}


\noindent
Large foundational models like LLMs have demonstrated remarkable
accuracy across a wide range of tasks including
writing~\cite{DBLP:journals/corr/abs-2303-08774}, coding~\cite{DBLP:journals/corr/abs-2308-12950}, 
and many others~\cite{DBLP:journals/corr/abs-2112-09332,perf-agent,DBLP:journals/corr/abs-2504-05738,DBLP:journals/corr/abs-2502-01450}. 
Post-training these models with reinforcement learning (RL)
has become a key pillar in the training pipeline
to further enhance model capabilities 
in challenging reasoning tasks~\cite{DBLP:journals/corr/abs-2501-12948}.

Rollout is a key and performance-dominant phase in post-training (\textsection{\ref{sec:bg-motivation}}):
at each step,
the system feeds the model with a batch of prompts representing
target problems (e.g., math or coding problems) to be solved.
The model then generates tokens for each prompt in an attempt
to solve them.
When rollout finishes, the generated tokens are evaluated
and the model is updated accordingly.
Rollout naturally parallelizes the batch across multiple workers. 


\stitle{Problem statement and current solutions (\textsection{\ref{sec:bg-motivation}}). \,}
Rollout accounts for {75--80}\% of the post-training time.
This lengthy and inefficient execution is primarily due to 
long tails in generation and is further exacerbated by idle GPU time, 
caused by waiting for the slowest worker to finish.
The long tail stems from difficulty variation, with some problems requiring more tokens than others to solve.
Therefore, one worker may finish its assigned prompts quickly,
while another takes much longer to complete.
Importantly, waiting for all requests before updating the model
is necessary to ensure rapid and stable training convergence 
~\cite{DBLP:journals/corr/SchulmanWDRK17,DBLP:journals/corr/abs-2402-03300,dapo}. 

Given this hard constraint, we ask a specific but important question:
\emph{how to accelerate rollout in LLM post-training {without changing
how the model is trained}?}


Existing \emph{algorithm-agnostic} solutions fall into two categories,
neither of which adequately addresses the issue 
given the recent trend of ever-growing generation length as
LLMs are trained to solve more complex problems~\cite{swebench,agentbench,deepseek-r1}.
The first category is \emph{overlapping}, which utilizes idle GPUs in the rollout for other
phases in the training pipeline~\cite{rlhfuse}.
% It is ineffective because it does not directly speed up rollout:
% In common post-training with long-tailed generation,
% the speedup is only 1.1\,$\times$ on average 
% because rollout dominates the training time. 
It has been proven effective in short-context RL training but 
less effective in long-tailed training 
because it does not accelerate rollout directly, yielding 
1.1\,$\times$ average speedup.
The second category leverages \emph{parallelism},
\ie{utilizing idle or over-provisioned GPUs} to accelerate post-training~\cite{rlboost}.
Unfortunately,
the speedup is still limited (e.g., up to 1.2\,$\times$ with doubled GPUs)
because LLM generation is memory-bound and thus cannot fully utilize the extra computation power. 
%Worse still,
%GPU over-provisioning is not always feasible in practice ({\fig{fig:motiv-batchsize}}).}      

\stitle{This work: lossless and efficient speculative rollout. \,}We present {\sys}, a fast rollout system that retrofits speculative decoding~\cite{DBLP:conf/osdi/BehrensCSBKZ20,DBLP:conf/osdi/ChangG99,spec-action,spec-edge}---a
common technique for accelerating LLM inference---to LLM post-training. 
%\RX{Speculation is a classic and widely-adopted system technique 
%~\cite{DBLP:conf/osdi/BehrensCSBKZ20,DBLP:conf/osdi/ChangG99,spec-action,spec-edge}.}
Specifically, speculative decoding employs a fast path for sequential generation (decoding):
First, a ``draft'' model much smaller than the trained one 
generates a sequence of $n$ draft tokens.
Then, the trained model verifies these tokens. 
Once accepted, these tokens are treated as generated by the trained model.
The rationale is that
verifying $n$ tokens is {typically faster than generating them}---although 
using the same model---because
{verification processes multiple tokens in parallel.}

It is important to note that speculative decoding can provide lossless generation 
as the original decoding via exact token matching~\cite{DBLP:conf/acl/XiaYDW00L0S24}.
Moreover, the drafters are available at post-training~\cite{spd,specinfer}:
pre-training pipeline generates both small and large models,
following established scaling practices~\cite{qwen2.5,qwen3}.
Besides, n-gram-based drafters are model-free~\cite{pld,DBLP:conf/acl/HuWZZLCZ25}. 



While intuitive and effective in inference, applying vanilla speculative decoding
to rollout yields only {10\,\%} speedup at worst on common production configurations 
(\textsection{\ref{sec:eval-e2e}})
due to the following challenges. We further address them with 
two contributions.

\stitle{Computation-efficient rollout via decoupled speculation (\textsection{\ref{sec:design-decouple}}). \,}
First, the effectiveness of speculative decoding depends mainly on the verification overhead,
which is proportional to the per-worker batch size. 
However, on common training configurations, e.g., per-worker batch size 128, (see {\fig{fig:motiv-batchsize}} (a)),
the verification could become slower than generation 
due to reaching the GPU computation limit,
also noticed by recent works~\cite{eagle-3,turbospec}.
Although the batch size gradually decreases during rollout as prompts finish,
it remains relatively large (\eg{$\geq$ 64}) for a considerable portion of the rollout time (\eg{20\,\%}),
so we need efficient execution on batch configurations 
that are unfriendly to speculation.

Our key observation is that there exist opportunities to accelerate verification because
vanilla speculative execution limits the GPU time provisioned to (compute-hungry) verification
due to a coupled draft-$n$-then-verify dependency:
the verifier must wait for GPU-underutilized drafter jobs.
Such a dependency is necessary to ensure high speculation efficiency for \emph{all requests} 
(required in inference~\cite{turbospec})
because without it, 
the drafter could further waste all $[n,...)$ tokens if verification rejects one token in $[1,n]$.
However, such a dependency is overly conservative for rollout 
where only the batch completion time matters,
so trading off a few requests to accelerate others is beneficial.
Based on this, we propose decoupled speculation to relax the dependency 
between draft and verification. 
The drafter can continue drafting without waiting for verification to finish,
thereby allowing subsequent verifications to start earlier to maximize GPU time provisioned to verification.
The decoupling comes at the cost of decreased acceptance rate due to more in-flight unverified tokens.
Fortunately, many requests only have a slightly decreased acceptance rate,
so the computation gain offsets the acceptance loss.
Moreover, we further propose a
feedback-based draft window mechanism to dynamically minimize wasted tokens due to such aggressive drafting 
via online reconfiguration.



\stitle{Effective draft adaptation via Fastest-of-N speculation (\textsection{\ref{sec:design-bon}}). \,}
Second, speculative decoding has multiple candidate draft methods,
and it is a common practice to select 
a single drafting method for all requests based on profiling or experience~\cite{spd,specinfer}.
However, we found in rollout that different requests suit different draft methods, 
so using one fixed method might lead to inefficient speculation 
for stragglers. 
However,
it is challenging to select the optimal one for all requests because
for each request, the optimal method not only depends on how fast the drafter is,
but also on the acceptance rate of verification.
The latter is unknown in advance as it depends on the model output. 
%Thus, the uniformly applied darting method inevitably leads to suboptimal speedup for 
%some (tailed) requests.
Meanwhile, it is impractical to run multiple draft methods for each request
as it requires proportional extra verifiers and thus extra GPUs.

We propose Fastest-of-N speculation that approximates the 
best draft methods for all requests with the following three techniques. 
First, we propose a new abstraction called \emph{draft ladder} that 
establishes a mapping from various acceptance rates to speedup given a pool of draft methods,
which can be constructed offline. 
Given the ladder, at the beginning of rollout,
we select one estimated fastest {draft method}
using the average acceptance rate of all requests profiled so far.
The rationale is that even though each requests' acceptance rate varies significantly,
the average acceptance rate over a (large) batch of requests is statistically stable~\cite{hoeffding1963probability}.
Thus, the selected draft method is likely to be close to the optimal for the batch. 
Finally, when the batch size decreases during rollout and
the initial selection becomes suboptimal,
we utilize the freed GPUs from finished workers to launch more draft methods in parallel for long-tailed requests,
and the request finishes when the fastest draft method 
generates the end token (EOS) that is accepted by the verifier.
Running multiple draft methods in parallel increases
the likelihood of selecting the fastest method
with unknown acceptance rates, while utilizing the idle GPUs
avoids extra GPU usage.


\stitle{Demonstration. }
{\sys} preserves the exact rollout process for any training algorithm,
so algorithm designers can seamlessly use it without worrying about potential accuracy loss.
We built {\sys} on {veRL}~\cite{hybridflow}, a state-of-the-art
post-training framework. 
% {\sys} is easy to integrate into other training frameworks 
% because it can serve as a clean drop-in replacement for the inference part 
% in any post-training framework, unlike other works that require modifying both the training and inference engines.
{\sys} is easy to integrate into other training frameworks, as it serves as a drop-in replacement for the inference component 
of any post-training pipeline, without requiring changes to the training logic.
Extensive evaluations show that on various real-world training traces 
and common post-training algorithms including 
{GRPO~\cite{DBLP:journals/corr/abs-2402-03300,deepseek-r1}, DAPO~\cite{dapo}, and PPO~\cite{DBLP:journals/corr/SchulmanWDRK17}} , 
compared to the state-of-the-art solutions like {veRL and RLHFuse},
{\sys} accelerates the end-to-end training time by {1.4--2.3}\,$\times$.
Also, we are up to {2.6}\,$\times$
{faster than vanilla speculative rollout baselines in different traces.}